\documentclass[11pt,a4paper]{article}
\usepackage[utf8]{inputenc}
\usepackage[T1]{fontenc}
\usepackage[french]{babel}
\usepackage[margin=2.4cm]{geometry}
\usepackage{amsmath,amssymb}
\usepackage{booktabs}
\usepackage{xcolor}
\usepackage[colorlinks=true,linkcolor=blue!50!black,citecolor=blue!50!black]{hyperref}
\usepackage{tcolorbox}
\tcbuselibrary{skins,breakable}
\usepackage{titlesec}
\usepackage{enumitem}
\usepackage{parskip}
\setlength{\parskip}{0.4em}

\definecolor{bleuprincip}{HTML}{1F4E78}
\definecolor{bleuclair}{HTML}{2E74B5}
\definecolor{orangealerte}{HTML}{C55A11}
\definecolor{vertok}{HTML}{548235}

\titleformat{\section}{\Large\bfseries\color{bleuprincip}}{\thesection}{0.6em}{}
\titleformat{\subsection}{\large\bfseries\color{bleuclair}}{\thesubsection}{0.6em}{}

\newtcolorbox{intuition}[1][]{colback=bleuclair!7,colframe=bleuclair,boxrule=0.5pt,arc=2pt,
  left=6pt,right=6pt,top=5pt,bottom=5pt,breakable,title={\small\bfseries Intuition},#1}
\newtcolorbox{definitionbox}[1][]{colback=gray!7,colframe=black!60,boxrule=0.4pt,arc=1pt,
  left=6pt,right=6pt,top=5pt,bottom=5pt,breakable,#1}
\newtcolorbox{resultatbox}[1][]{colback=vertok!8,colframe=vertok,boxrule=1pt,arc=2pt,
  left=7pt,right=7pt,top=5pt,bottom=5pt,breakable,#1}
\newtcolorbox{alertebox}[1][]{colback=orangealerte!7,colframe=orangealerte,boxrule=0.6pt,arc=2pt,
  left=7pt,right=7pt,top=5pt,bottom=5pt,breakable,title={\small\bfseries Honnêteté sur le résultat},#1}

\newcommand{\cit}[1]{\textsuperscript{\hyperlink{ref#1}{[#1]}}}

\title{\vspace{-1.3cm}\textcolor{bleuprincip}{\Huge\bfseries GDPNow-Maroc}\\[0.3cm]
\textcolor{black}{\Large Test hors-échantillon --- Comparaison directe avec Higgins (2014)}\\[0.15cm]
\textcolor{gray}{\large Méthodologie, résultats et honnêteté sur les limites}}
\author{}
\date{\today}

\begin{document}
\maketitle
\thispagestyle{empty}

\begin{abstract}
\noindent Ce document soumet le modèle complet GDPNow-Maroc à un \textbf{vrai test hors-échantillon} : les coefficients déjà estimés (jamais retouchés) sont appliqués à 20 trimestres --- 2021-T2 à 2026-T1 --- qu'aucune étape du pipeline n'a vus pendant l'estimation. Le résultat (RMSFE = 4,73 points de croissance annualisée) est comparé directement à celui de Higgins\cit{1} (RMSFE = 1,15, sur 14 ans de vraies prévisions temps réel).
\end{abstract}

\tableofcontents
\vspace{0.5cm}

% ============================================================================
\section{Pourquoi ce test, et pourquoi il est différent de tout ce qui précède}
% ============================================================================

\begin{intuition}
Toutes les corrélations rapportées jusqu'ici dans ce travail (0,671 à l'Étape 6, par exemple) étaient calculées \textbf{en-échantillon} --- sur les mêmes données qui ont servi à estimer les coefficients. Un modèle peut sembler bien ajusté simplement parce qu'il a « appris » le bruit spécifique de ces données, sans que ça garantisse rien sur sa capacité à prévoir un vrai trimestre à venir. Higgins ne rapporte jamais de résultat en-échantillon dans son évaluation finale --- uniquement du hors-échantillon (Table 5 de son article).
\end{intuition}

% ============================================================================
\section{Méthode --- pourquoi aucune réestimation n'était nécessaire}
% ============================================================================

\begin{definitionbox}
Le train (T1-2010--T1-2021) a été fixé une fois pour toutes, dès la Phase 2, et jamais dépassé pour \textbf{estimer} quoi que ce soit. Mais nos données réelles (VA, indicateurs) couvrent en réalité jusqu'à 2026 --- les trimestres 2021-T2 à 2026-T1 (20 trimestres) sont donc \textbf{déjà disponibles}, sans avoir jamais influencé la moindre estimation.
\end{definitionbox}

\begin{equation}
\underbrace{\text{Coefficients figés (train 2010--2021)}}_{\text{jamais réestimés}} \ \xrightarrow{\text{appliqués à}}\ \underbrace{\text{2021T2--2026T1}}_{\text{jamais vu}} \ \xrightarrow{\text{comparé à}}\ \underbrace{\text{le réalisé}}_{\text{déjà connu}}
\end{equation}

\subsection{Ce qui est réutilisé sans modification}

Les coefficients du BVAR (Phase 1, 1296 coefficients), les 12 bridge equations (avec leur $q,r$ déjà choisis par AIC), les 4 modèles AR (avec leur $p^\star$ déjà choisi), et les 16 $\delta$ (bornés, Étape 5) --- \textbf{aucun n'est retouché}. Seule la fenêtre de dates sur laquelle on les \emph{applique} s'étend.

% ============================================================================
\section{Résultat}
% ============================================================================

\begin{resultatbox}
\begin{center}
\begin{tabular}{lr}
\toprule
\textbf{Indicateur} & \textbf{Valeur} \\
\midrule
$n$ (trimestres jamais vus) & 20 \\
Corrélation (prévu, réalisé) & 0,541 \\
RMSFE (croissance annualisée) & 4,73 points \\
\bottomrule
\end{tabular}
\end{center}
\end{resultatbox}

% ============================================================================
\section{Comparaison directe avec Higgins (2014), Table 5}
% ============================================================================

\begin{center}
\small
\begin{tabular}{lrl}
\toprule
\textbf{Modèle} & \textbf{RMSFE} & \textbf{Contexte} \\
\midrule
AR(2) roulant (Higgins) & 2,43 & USA, 2000--2013, vrai temps réel \\
BVAR trimestriel (Higgins) & 2,28 & USA, 2000--2013, vrai temps réel \\
Modèle de suivi mensuel (Higgins) & 1,22 & USA, 2000--2013, vrai temps réel \\
\textbf{GDPNow complet (Higgins)} & \textbf{1,15} & USA, 2000--2013, vrai temps réel \\
\midrule
\textbf{GDPNow-Maroc (ce travail)} & \textbf{4,73} & Maroc, 2021--2026, pseudo-temps réel \\
\bottomrule
\end{tabular}
\end{center}

\begin{alertebox}
\textbf{Notre RMSFE est environ 4 fois plus élevé que celui du modèle complet de Higgins} --- et même supérieur à son AR(2) roulant le plus simple (2,43). C'est un résultat honnête à assumer, pas à minimiser. Trois différences majeures, structurelles, expliquent en grande partie cet écart :
\begin{enumerate}[leftmargin=1.4em]
\item \textbf{$n=20$ contre $n\approx56$} pour Higgins --- notre estimation du RMSFE repose sur un échantillon près de 3 fois plus petit, donc bien moins stable statistiquement.
\item \textbf{Densité et qualité des données} --- Higgins dispose de plus de 200 séries mensuelles américaines, avec des décennies d'historique fiable ; notre travail a dû composer, branche par branche, avec des indicateurs souvent faibles ou en AR pur (7 branches sur 16, faute de lien statistique exploitable).
\item \textbf{Une seule fenêtre de test, pas une validation glissante} --- le RMSFE de Higgins moyenne 56 réestimations indépendantes (Table 5), la nôtre n'en fait qu'une seule, jamais répétée.
\end{enumerate}
\end{alertebox}

% ============================================================================
\section{Ce que ce résultat permet de conclure, honnêtement}
% ============================================================================

\begin{resultatbox}
\textbf{Le modèle capte un vrai signal} --- une corrélation de 0,541 sur des données jamais vues n'est pas un hasard statistique pur, c'est un lien réel, quoique modeste. \textbf{Mais il est loin d'atteindre la précision du modèle original de Higgins}, pour des raisons largement structurelles (richesse des données disponibles, taille de l'échantillon de test) plutôt qu'une erreur de méthode identifiable et corrigible.
\end{resultatbox}

\section*{Références}
\addcontentsline{toc}{section}{Références}
\begin{enumerate}[leftmargin=1.6em, label=\textbf{[\arabic*]}]
\item \hypertarget{ref1}{} Higgins, P. (2014). \textit{GDPNow: A Model for GDP ``Nowcasting''}. Federal Reserve Bank of Atlanta, Working Paper 2014-7, Table 5.
\end{enumerate}

\end{document}
